\documentclass[11pt]{article}
\usepackage{amsmath,amssymb,amsthm,mathtools}
\usepackage[margin=1in]{geometry}

\newtheorem{theorem}{Theorem}
\newtheorem{definition}{Definition}
\newtheorem{lemma}{Lemma}

\title{Human-Readable Proof Outline for \texttt{kasami\_is\_ab}}
\author{RequestProject / Kasami AB}
\date{}

\begin{document}
\maketitle

\section{What the theorem says}
The Lean theorem \texttt{kasami\_is\_ab} proves that the Kasami power function
\[
  f(x) = x^{d(k)}, \qquad d(k) = 2^{2k} - 2^k + 1,
\]
is \emph{Almost Bent} on a finite field $F$ of characteristic $2$, provided that:
\begin{itemize}
  \item $|F| = 2^n$,
  \item $n \ge 1$ and $n$ is odd,
  \item $k \ge 1$,
  \item $\gcd(k,n) = 1$.
\end{itemize}

In the formal development, this is expressed as
\[
  \mathrm{IsAB}(hcard, f),
\]
where \texttt{hcard} is the hypothesis $|F| = 2^n$ and $f(x) = x^{d(k)}$.

\section{Main proof steps}
The Lean proof follows a short reduction chain:
\begin{enumerate}
  \item Reduce the AB goal to a general moment criterion \texttt{ab\_from\_moments}.
  \item Supply the Parseval identity for the Walsh spectrum.
  \item Supply the fourth-moment identity coming from APN.
  \item Supply the divisibility of every Walsh coefficient by a power of $2$.
  \item Conclude that the Walsh spectrum has the AB shape, so the function is Almost Bent.
\end{enumerate}

\section{Step 1 in detail: reduce to the general AB criterion}
The theorem is not proved directly from scratch. Instead, the function is inserted into a general theorem from \texttt{Walsh/AB.lean}:
\[
  \texttt{ab\_from\_moments}.
\]
That theorem says that if a function $f : F \to F$ satisfies three conditions for every nonzero shift $a$:
\begin{itemize}
  \item a Parseval-type second moment identity,
  \item a fourth moment identity,
  \item a divisibility statement for its Walsh coefficients,
\end{itemize}
then $f$ is AB.

So the role of \texttt{kasami\_is\_ab} is to verify those three inputs for the specific Kasami monomial.

Concretely, the proof first fixes the field hypotheses and the Kasami parameters $n$ and $k$. Then it applies \texttt{ab\_from\_moments} to the function
\[
  f(x) = x^{d(k)}.
\]
At that point, the only remaining job is to fill in the three moment/divisibility assumptions using earlier Kasami-specific lemmas.

\section{Definitions used in the theorem}
This section explains the formal objects that appear in the proof.

\subsection{Kasami exponent}
The exponent is
\[
  d(k) = 2^{2k} - 2^k + 1.
\]
The function under study is the power map
\[
  x \mapsto x^{d(k)}.
\]
This is the standard Kasami monomial.

\subsection{Almost Bent (AB)}
In the Lean development, the Almost Bent predicate is defined as follows. For a function $f : F \to F$ and a size hypothesis $|F| = 2^n$,
\[
  \mathrm{IsAB}(f) :\!\iff \forall a \in F,\ a \neq 0 \rightarrow \forall b \in F,\
  \mathrm{walsh}(f,a,b)^2 = 0 \ \text{or}\ \mathrm{walsh}(f,a,b)^2 = 2^{n+1}.
\]
Equivalently, for every nonzero shift $a$, every Walsh coefficient squared is forced to be either $0$ or the fixed value $2^{n+1}$.

This is the precise AB magnitude pattern used by the theorem \texttt{ab\_from\_moments}.

\subsection{Walsh transform}
The Walsh coefficient of a function $f : F \to F$ at $(a,b)$ is defined by
\[
  \mathrm{walsh}(f,a,b) = \sum_{x \in F} \chi\bigl(a x + b f(x)\bigr),
\]
where the sign character is
\[
  \chi(z) = \begin{cases}
  1 & \text{if } \mathrm{Tr}(z)=0,\\
  -1 & \text{if } \mathrm{Tr}(z)\neq 0.
  \end{cases}
\]
Here $\mathrm{Tr} : F \to \mathbf{GF}(2)$ is the absolute trace.

The proof uses three facts about these coefficients:
\begin{itemize}
  \item Parseval gives a second-moment identity.
  \item The APN property gives a fourth-moment identity.
  \item A separate divisibility theorem shows each coefficient is divisible by a power of $2$.
\end{itemize}

\subsection{Walsh divisibility framework}
The Walsh divisibility argument is the step that turns the Kasami exponent identity into the exact power-of-two divisibility needed by the AB theorem.

The chain of reasoning is:
\begin{enumerate}
  \item Start with the Kasami Walsh coefficient $\mathrm{walsh}(x \mapsto x^{d(k)}, a, b)$.
  \item Use the bijection $x = y^{2^k+1}$ to rewrite the sum as a character sum of a quadratic function.
  \item Observe that the resulting exponent pattern is quadratic in the finite-field sense, because its third discrete derivative vanishes.
  \item Apply the general quadratic Gauss-sum divisibility theorem: over a field of size $2^n$ with $n$ odd, such a sum is divisible by $2^{(n+1)/2}$.
\end{enumerate}

The relevant formal lemmas are:
\begin{itemize}
  \item \texttt{gold\_pow\_bijective}: the Gold map $y \mapsto y^{2^k+1}$ is a permutation under the Kasami hypotheses.
  \item \texttt{walsh\_kasami\_eq\_quadratic}: the Kasami Walsh sum is rewritten using that substitution.
  \item \texttt{sum\_gold\_third\_deriv\_zero}: the substituted polynomial has vanishing third derivative.
  \item \texttt{quadratic\_gauss\_sum\_div}: every quadratic character sum over $GF(2^n)$ with $n$ odd is divisible by $2^{(n+1)/2}$.
  \item \texttt{kasami\_walsh\_div}: the final Kasami divisibility statement.
\end{itemize}

Conceptually, this framework is a bridge from algebraic structure to arithmetic output: the exponent $d(k)$ is special enough that the Walsh sum becomes quadratic after a change of variables, and odd-dimensional quadratic sums are forced to have large $2$-adic valuation.

\subsection{Quadratic Gauss-sum divisibility}
The general theorem used at the end of the Walsh divisibility framework is \texttt{quadratic\_gauss\_sum\_div}. It says that if $Q : F \to F$ is quadratic in the sense that its third discrete derivative vanishes, and if $|F|=2^n$ with $n$ odd, then the character sum
\[
  \sum_{x \in F} \chi\bigl(Q(x)\bigr)
\]
is divisible by
\[
  2^{(n+1)/2}.
\]

The proof has two stages. First, one squares the sum and rewrites
\[
  \left(\sum_x \chi(Q(x))\right)^2
\]
as a double sum of the form
\[
  \sum_u \sum_y \chi(Q(u+y)+Q(y)).
\]
For fixed $u$, this inner expression is split into a constant factor and an additive character sum involving
\[
  B(u,y)=Q(u+y)+Q(u)+Q(y)+Q(0).
\]
The quadratic hypothesis implies that $B(u,\cdot)$ is additive, so the additive-character lemma forces its sum to be divisible by $|F|=2^n$. That gives $2^n \mid S^2$ for the original sum $S$.

Second, since $n$ is odd, a $2$-adic parity lemma upgrades divisibility of the square to divisibility of the sum itself: from $2^n \mid S^2$ it follows that
\[
  2^{(n+1)/2} \mid S.
\]

This is the exact mechanism behind the theorem \texttt{quadratic\_gauss\_sum\_div}, and it is the generic result that \texttt{kasami\_walsh\_div} instantiates after rewriting the Kasami Walsh sum into quadratic form.

\subsection{APN}
APN means \emph{almost perfect nonlinear}. For a function in characteristic $2$, this is the differential uniformity condition that every nonzero derivative equation has at most two solutions.
The Kasami APN theorem is proved elsewhere in the project and is imported here through the lemma \texttt{kasami\_is\_apn\_pred}.

\subsection{Bijectivity}
The proof also needs the Kasami map to be a permutation of the finite field. This is established by proving injectivity first and then using the finite-set fact that injective self-maps are bijective.

\subsection{The field hypotheses}
The conditions
\[
  |F| = 2^n, \qquad n \text{ odd}, \qquad n \ge 1
\]
ensure that the Walsh-spectrum theorem applies in the intended odd-dimensional characteristic-2 setting. The conditions on $k$ are the usual Kasami arithmetic assumptions: $k \ge 1$ and $\gcd(k,n)=1$.

\section{How the proof closes}
After Step 1 reduces the theorem to the general moment criterion, the remaining three obligations are supplied by previously proved lemmas:
\begin{itemize}
  \item Parseval is obtained from the bijectivity of the Kasami map.
  \item The fourth moment is obtained from the APN property together with bijectivity.
  \item Walsh divisibility is obtained from the dedicated Kasami Walsh-divisibility theorem.
\end{itemize}

Once those are in place, \texttt{ab\_from\_moments} yields that the Kasami monomial is AB.

\section{One-paragraph summary}
The theorem \texttt{kasami\_is\_ab} proves that the Kasami power map $x \mapsto x^{2^{2k} - 2^k + 1}$ is Almost Bent over a finite field of size $2^n$ when $n$ is odd and $k$ satisfies the standard Kasami coprimality assumptions. The proof is modular: instead of analyzing the Walsh spectrum directly, it reduces the AB claim to a general moment criterion, then fills in the Parseval identity, the fourth-moment identity, and a divisibility bound using earlier Kasami-specific results.

\end{document}
